% Generated draft: requires mathematical and source review.
\documentclass[11pt]{article}
\usepackage[margin=25mm]{geometry}
\usepackage{fontspec}
\setmainfont{Latin Modern Roman}
\usepackage{amsmath,amssymb,mathtools,unicode-math}
\setmathfont{Latin Modern Math}
\usepackage{enumitem,needspace,xurl,hyperref}
\setlength{\parindent}{0pt}
\setlength{\parskip}{5pt}
\setlength{\emergencystretch}{3em}
\newcommand{\N}{\mathbb{N}}
\newcommand{\Z}{\mathbb{Z}}
\newcommand{\Q}{\mathbb{Q}}
\newcommand{\R}{\mathbb{R}}
\newcommand{\C}{\mathbb{C}}
\newcommand{\F}{\mathbb{F}}
\newcommand{\A}{\mathbb{A}}
\newcommand{\PP}{\mathbb{P}}
\newcommand{\E}{\mathbb{E}}
\newcommand{\eps}{\varepsilon}
\newcommand{\dd}{\,\mathrm{d}}
\begin{document}
\section*{UnsolvedMath ID: 30001649}
\textbf{Source problem code:} catalogue entry
\subsection{Definitions and mathematical statement}
An oval of a real planar Hamiltonian system is understood here, conditionally on the usual real interpretation, as a smooth simple closed component \(\gamma_h\) of a regular level \(H^{-1}(h)\) of a Hamiltonian \(H\). Regularity means \(\nabla H\ne 0\) on \(\gamma_h\). Write \(J\subseteq\R\) for the unspecified parameter interval of the supplied family \(\{\gamma_h\}_{h\in J}\). The real domain, regularity of \(H\), and interval \(J\) are not explicitly stated in the record; this interpretation must therefore be confirmed in mathematical review.

For the polynomial functions \(f\) and \(g\) appearing in the source, define the Abelian integral
\[
I(h)=\oint_{\gamma_h}\bigl(f(x,y)\,\mathrm{d}y-g(x,y)\,\mathrm{d}x\bigr),\qquad h\in J.
\]
The source does not specify the orientation of the ovals. Choosing a coherent orientation defines the displayed integrals; reversing that orientation changes their signs.

\textbf{Source question.} For a polynomial perturbation, depending on a small parameter \(\varepsilon\), of the planar Hamiltonian system having these ovals, determine how many of the unperturbed periodic orbits persist or give rise to periodic orbits of the perturbed system as \(\varepsilon\to0\).

\textbf{Unresolved qualifications.} The record does not give the perturbed differential equations or identify precisely how \(f\) and \(g\) enter them. It supplies no degree bounds, explicit Hamiltonian, interval endpoints, or assumptions on degeneracies of \(I\). It also does not specify whether the desired count concerns all periodic orbits or only limit cycles, meaning isolated periodic orbits; whether zeros or cycles are counted with multiplicity; whether the target is an exact count for fixed data or a uniform bound over a class; or whether bifurcations at boundary levels are included. Consequently the supplied question does not determine a unique precise counting problem. No equality between the requested number and a zero count of \(I\) is asserted here, and the case \(I\equiv0\) is not excluded by the source.
\subsection{Short English statement}
The problem asks how many periodic trajectories of a planar Hamiltonian system survive or generate periodic trajectories under a small polynomial perturbation, with the displayed Abelian integral as associated data. The perturbation equations, counting convention, degree restrictions, and treatment of degenerate and boundary cases require clarification. The supplied open designation is dated curation, not an independently verified mathematical conclusion.
\subsection{Sources}
\begin{itemize}
\item Local registry: \texttt{data/unsolvedmath/problems.json} and supplementary registry \texttt{catalog/problems.json}; supplied complete source record for permanent UnsolvedMath ID \(30001649\), titled \emph{Limit Cycles and Zeros of Abelian Integrals}, in \emph{Oberwolfach Reports Open Problems}. Exact source problem code: \texttt{catalogue entry}. These local registry references identify the supplied record; their contents were not independently inspected.
\item Supplied report citation: \emph{Mini-Workshop: Algebraic and Analytic Techniques for Polynomial Vector Fields} (2011), \emph{Oberwolfach Reports}, pp.\ 3311--3348, DOI \texttt{10.4171/owr/2010/57}. \url{https://doi.org/10.4171/owr/2010/57}. The record gives proposed year 2010 and names no individual proposer. The differing report-title year and proposed year are retained without reconciliation.
\item Supplied literature source, labelled ``original Oberwolfach problem report'': \url{https://doi.org/10.4171/owr/2010/57}.
\item Supplied literature source, labelled ``Abelian integrals and limit cycles'': \url{https://ems.press/content/serial-article-files/46318}. No author or additional bibliographic details are supplied for this link.
\item Literature checked at: \texttt{2026-08-21}. This is supplied dated curation, not independently verified. The supplied assessment says that Poincar\'e--Pontryagin theory reduces first-order persistence to zeros of Abelian integrals, while the general zero-count and limit-cycle problem remains open. The background explicitly requests independent verification of time-sensitive status and cited claims; the hypotheses and scope of this assessment are not provided.
\item The registry labels the problem \texttt{open} and its statement \texttt{corrected\_verified}, with a supplied claim of manual semantic, typographic, and LaTeX review. These are source metadata, not independent verification of mathematical completeness or current research status.
\end{itemize}
\end{document}
